\BourbakiExercises{Exercises}

\begin{enumerate}
\def\labelenumi{\arabic{enumi})}
\tightlist
\item
  Let \(D\) be a field.
\end{enumerate}

\begin{enumerate}
\def\labelenumi{\alph{enumi})}
\item
  Let E be a proper subfield of D whose multiplicative group \(E^{*}\) has finite index in \(D^{*}\) . Prove that D is finite (for a sequence \((a_{n})\) of distinct elements of E and an element x of D---E, consider the classes mod \(E^{*}\) of the elements \(x + a_{n}\) ).
\item
  Prove that an element of D that has only finitely many conjugates belongs to the center Z of D (apply a) to the commutant of this element in D, which is a field).
\item
  Deduce that a polynomial in Z{[}X{]} that has a root in D-Z has infinitely many roots in D-Z.
\end{enumerate}

¶ 2) Let A be a K-algebra and m an integer. Suppose that every element of A is algebraic of degree \(\leqslant m\) over K.

\begin{enumerate}
\def\labelenumi{\alph{enumi})}
\item
  Prove that if A is a primitive ring, then there exist an integer \(r \leqslant m\) and a field D whose elements all have degree \(\leqslant m\) over K such that A is isomorphic to the matrix algebra \(\mathbf{M}_{r}(\mathrm{D})\) . (If A is not simple, then deduce from the density theorem that for every integer n, there exist a subalgebra \(A_{n}\) of A and a surjective homomorphism \(A_{n} \to M_{n}(\mathrm{D})\) . Observe that the algebra \(\mathbf{M}_{n}(\mathrm{D})\) contains elements X such that \(X^{n-1} \neq 0\) and \(X^{n} = 0\) .)
\item
  Suppose that the field K is infinite. Prove that the algebra \(A/\Re(A)\) is semisimple and admits at most m simple components (observe that the algebra \(K^{n}\) contains elements of degree n over K; deduce that A admits at most m maximal two-sided ideals, and conclude using a)). Prove that if K is, moreover, perfect, then \(A/\Re(A)\) has finite rank over K.
\end{enumerate}

¶ 3) Let K be a commutative field. We say that a K-algebra A is algebraic if all its elements are algebraic over K. Every algebra of finite rank is algebraic; every subalgebra and quotient algebra of an algebraic algebra is algebraic. The radical of an algebraic K-algebra is a nil ideal (VIII, p.~166, Exercise 5).

\begin{enumerate}
\def\labelenumi{\alph{enumi})}
\item
  Prove that an algebraic algebra is a pseudoregular ring (VIII, p.~178, Exercise 4; observe that for every \(x \in \mathrm{A}\) , there exist an element \(y \in \mathrm{K}[x]\) and an integer \(k\) such that \(x^k = x^{k+1}y\) ).
\item
  Let A be an algebraic K-algebra without any nilpotent elements other than 0. Prove that A is isomorphic to a subalgebra of a product of fields. (Observe that every quotient algebra of A has the same properties. Prove that if A is primitive, then it is a field, by reasoning as in Exercise 2, a). In the general case, use Exercise 12 of VIII, p.~168.)
\end{enumerate}

¶ 4) Let K be a finite field and A an algebraic K-algebra (Exercise 3).

\begin{enumerate}
\def\labelenumi{\alph{enumi})}
\item
  Suppose that A is a field. Prove that A is commutative. (Let \(x\) be a noncentral element of A. Consider the field \(Z(x)\) , where \(Z\) is the center of A, and prove that there exists an element \(y \in A\) such that \(yxy^{-1} = x^r\) , with \(x^r \neq x\) . Deduce a contradiction by considering the subfield \(K(x,y)\) of A.)
\item
  Suppose that A does not contain any nonzero nilpotent element. Prove that the algebra A is commutative (use a) and Exercise 3).
\end{enumerate}

\begin{enumerate}
\def\labelenumi{\arabic{enumi})}
\setcounter{enumi}{4}
\item
  Let K be a finite field, m an integer, and A a K-algebra without radical whose elements all have degree \(\leqslant m\) . Prove that there exists a finite extension L of K such that A is isomorphic to a subalgebra of a product \((\mathbf{M}_{r}(\mathrm{L}))^{\mathrm{I}}\) whose projections are simple algebras (cf.~Exercise 2, a)). Also prove the converse.
\item
  Let A be a ring such that for every \(a \in \mathrm{A}\) , there exists an integer \(n(a) > 1\) such that \(a^{n(a)} = a\) . Prove that A is isomorphic to a product of finitely many reduced commutative algebraic algebras over finite fields. (Observe that A is annihilated by an integer \(n\) , and then that for every prime number \(p\) , the \(p\) -component \(\mathrm{A}_p\) of the additive group A is a two-sided ideal annihilated by \(p\) . Apply Exercise 4, b).) Also prove the converse.
\item
  \begin{enumerate}
  \def\labelenumii{\alph{enumii})}
  \tightlist
  \item
    Prove that if the field K has property \((\mathrm{C}_{1})\) (VIII, p.~357), then so does every algebraic extension of K. (Reduce to the case of an extension L of K of finite degree r. Prove that if F is a homogeneous polynomial function of degree d on an L-vector space V, then the function \(N_{L/K} \circ F\) is homogeneous polynomial of degree rd on the K-vector space V.)
  \end{enumerate}
\end{enumerate}

\(^{*}b)\) Assume from now on that K is algebraically closed. Let \(F_{1},\ldots,F_{n}\) be homogeneous polynomials of degree \textgreater0 in m variables. Prove that if n \textless{} m, then there exists an \(a \neq 0\) in \(K^{n}\) such that \(F_{1}(a) = \cdots = F_{n}(a) = 0\) (use AC, VIII, §2, n° 4, p.~19, corollaire 1 du théorème 3 and AC, VIII, §3, n° 1, p.~24, proposition 2). c) Deduce from b) that the field \(K(X)\) has property \((C_{1})\) .

(Let f be a homogeneous polynomial of degree d in n variables whose coefficients are polynomial in X, and let \(\nu\) be the maximum of the degrees of these coefficients. Let N be an integer, and let \(G_{1},\ldots,G_{n}\) be polynomials in X of degree \(\leqslant N\) . Prove that the equality \(\mathrm{F}(\mathrm{G}_{1}(\mathrm{X}),\ldots,\mathrm{G}_{n}(\mathrm{X}))=0\) is equivalent to the annihilation of \(Nd+\nu+1\) homogeneous polynomials of degree \textgreater0 in the \(n(N+1)\) coefficients of the \(G_{i}\) .)

\(d\) ) Deduce from \(a\) ) and \(c\) ) that an extension of an algebraically closed field of transcendence degree \(\leqslant 1\) has property \((\mathrm{C}_1)\) (``Tsen's theorem'').*

